\hspace{3mm}

\centerline{\bf \Huge Группы, кольца, поля.}

\begin{flushright}
    \scriptsize{И я спросил его, \href{https://www.youtube.com/watch?v=xwvKcWa2vLM}{а в каком кольце}?\\На что доцент ответил  " а что такое кольцо?"\ }
\end{flushright}

Множество $M$ называется \textit{замкнтуым относительно операции $*$}, если для любых элементов $a$ и $b$ из множества $M$ выполнено, что $a*b$ тоже принадлежит множеству $M$

\textit{Упражнение.} Придумайте какое-нибудь множество и операцию, относительно которой множество не будет замкнутым. (Не путайте с тем, что не у всех есть обратный элемент. Например, целые числа замкнуты относительно умножения, но не у всех есть обратный относительно умножения.)

\textbf{Группа} - это множество $G$ с заданной на нем бинарной операцией $*$ (если говорится, что задана бинарная операция, то уже заранее подразумевается, что множество замкнуто относительно этой операции) и в котором выполнены три аксиомы:

1) Ассоциативность: $\forall a, b, c \in G \ \text{выполнено, что} (a*b)*c = a*(b*c)$

2) Существование нейтрального элемента: $\exists e \in G$ такой, что $\forall a \in G \ a*e = e*a = a$

3) Существование обратного элемента: $\forall a \in G \ \exists b $ такой, что $a*b=b*a = e$. Элемент $b$ называется \textbf{обратным к a} и обозначается $a^{-1}$

\textit{Упражнение.} 

$a$) Покажите, что $\mathbb{Z}, \mathbb{Z}_n, \mathbb{Z}_p, \mathbb{Q}$ - являются группами и укажите относительно каких операций. 

$b$) Придумайте какое-нибудь множество, которые было бы группой сразу для двух операций. 

$c$) Придумайте группы, которая состаяла бы не из чисел.

Группы, в которых операция коммутативна (то есть $a*b=b*a$) называются \textit{абелевыми группами}, в честь норвежского математика Нильса Хенрика Абеля.

\textbf{Кольцо} — множество $R$, на котором заданы две бинарные операции:$+, \times$ (называемые сложение и умножение), со следующими свойствами, выполняющимися для любых $a, b, c \in R:$

1. $a+b=b+a$ — коммутативность сложения;

2. $a+(b+c) = (a+b)+c$ - ассоциативность сложения;

3. $\exists 0 \in R : a+0 = 0+a = a$ - существование нейтрального элемента относительно сложения.

\newpage

\hspace{3mm}

4. $\forall a \in R \ \exists b \in R: a+b = b+a = 0$ - существование противоположного (обратного) элемента относительно сложения;

5. $(a\times b)\times c = a\times (b \times c)$ - ассоциативность умножения;

6.
$
\begin{cases}
    a\times(b+c) = a\times b + a\times c  - \text{дистрибутивность умножения относительно сложения}\\
    (b+c)\times a = b\times a + b\times c - \text{дистрибутиность сложения относительно умножения}
\end{cases}
$

\hspace{5mm}

\textit{Упражнение.} Покажите, что $\mathbb{Z}$ - целые числа являются кольцом. Приведите пример ещё хотя бы 5 колец.

Как можно видеть кольцо является абелевой группой относительно сложения, а вот относительно умножения есть ограничения. Если убрать эти ограничения, то мы получим определение поля. 

\textbf{Поле} — множество $\mathbb{F}$, на котором заданы две бинарные операции:$+, \times$ (называемые сложение и умножение), со свойствами абелевой группы относительно операции $+$, абелевой группы относительно операции $*$ и со всеми законами дистрибутивности. 

В поле всегда нейтральный элемент по сложению и умножению разные. Иначе говоря 0 не равен 1. 

\textit{Упражнение.} Приведите примеры конечных и бесконечных полей.


\hspace{3mm}

\centerline{\textbf{\Large Задачи для самопроверки.}}

\hspace{3mm}

Укажите, является ли данное множество кольцом / полем? :

1) $\mathbb{Z}, \mathbb{Q}, \mathbb{R}, \mathbb{R}\times \mathbb{R}, \mathbb{Z}_n, \mathbb{Z}_p$

2) Множество многочленов с коэффициентами из какого-то поля $\mathbb{F}$.

3) Множество комплексных чисел $\mathbb{C}$, то есть чисел вида $a + bi$, где $a, b \in \mathbb{R}$, а элемент $i$ это \textit{мнимая единица} для которой верно $i^2 = -1$ 

4) Множество \textit{целых Гауссовых чисел} то есть числа вида $a+bi$, где $a, b \in \mathbb{Z}$